\documentclass[11pt,a4paper]{article}
\usepackage[margin=2.4cm]{geometry}
\usepackage[T1]{fontenc}
\usepackage[utf8]{inputenc}
\usepackage{amsmath,amssymb}
\usepackage{booktabs}
\usepackage{xcolor}
\usepackage[numbers,sort&compress]{natbib}
\usepackage[colorlinks=true,linkcolor=blue!55!black,citecolor=green!45!black]{hyperref}
\usepackage{enumitem}

\newcommand{\GHI}{\ensuremath{\Gamma_{\rm HI}}}
\newcommand{\GHeI}{\ensuremath{\Gamma_{\rm HeI}}}
\newcommand{\GHeII}{\ensuremath{\Gamma_{\rm HeII}}}
\newcommand{\nion}{\ensuremath{\dot{n}_{\rm ion}}}
\newcommand{\cMpc}{\ensuremath{{\rm cMpc}}}
\newcommand{\flag}[1]{\textcolor{red}{\textsf{[#1]}}}

\title{\bf Photoionization rates per target atom at the end of reionization,\\
in units directly comparable to a cosmic-ray ionization rate}
\author{Prepared for L.\ Carvalho}
\date{9 September 2026}

\begin{document}
\maketitle

\begin{abstract}
\noindent
Converts the ionizing output of the reionization-era source population into
$\Gamma_i$, the photoionization rate \emph{per target particle} in s$^{-1}$,
for H\,\textsc{i}, He\,\textsc{i} and He\,\textsc{ii} --- the same units as a
cosmic-ray ionization rate $\zeta$.  The hydrogen rate is taken from direct
measurement; the helium rates are derived as ratios to it, which cancels the
background normalisation and leaves only the spectral shape and the
photoionization cross-sections.  Three spectral bases are compared.
Script: \texttt{ionizing\_rates\_vs\_cosmic\_rays.py}.
\end{abstract}

%=====================================================================
\section{Assumptions}
%=====================================================================

\begin{enumerate}[leftmargin=1.7em]
\item \textbf{Redshift and phase.} All rates are quoted at $z=6.0$ and $z=4.9$,
      i.e.\ at and just after the end of reionization, where \GHI\ is measured.
      They apply to gas \emph{inside ionized regions}; see
      Section~\ref{sec:caveats}.
\item \textbf{Escaping, not intrinsic.} The background is the one that survives
      escape from galaxies and propagation through the IGM, because that is
      what \citet{Gaikwad2023} measure from the Ly$\alpha$ forest.
\item \textbf{Optically thin targets.} $\Gamma_i$ is computed for an
      unshielded atom bathed in the mean background.  This is a good
      approximation for H\,\textsc{i} and He\,\textsc{i} at $z\simeq6$ but
      \emph{not} for He\,\textsc{ii} (Section~\ref{sec:caveats}).
\item \textbf{Spectral shape.} Three bases, none of which is assumed correct
      a priori:
      \begin{itemize}[nosep]
      \item[(A)] a power law $J_\nu\propto\nu^{-\alpha_{\rm s}}$ with
            $\alpha_{\rm s}=2.0\pm0.6$, the effective ionizing-source index
            measured by \citet{Gaikwad2023};
      \item[(B)] the BPASS stellar SED of \citet{Eide2020}, Fig.~1, read off as
            a piecewise power law with indices $0.97$ ($13.6$--$24.6\,$eV),
            $2.18$ ($24.6$--$54.4\,$eV) and $23.8$ above the He\,\textsc{ii}
            edge.  \flag{Read off a published figure; I estimate the
            uncertainty at $\pm0.2$~dex per anchor point, which propagates to
            roughly $\pm0.3$ in each index.}
      \item[(C)] the Pop~III photon budgets of \citet{Murphy2021}, Table~4,
            used for photon-number ratios only.
      \end{itemize}
\item \textbf{Cross-sections} are the standard analytic fits of
      Verner et al.\ (1996) / Osterbrock \& Ferland (2006):
      $\sigma_{\rm HI}(\nu_0)=6.30\times10^{-18}$,
      $\sigma_{\rm HeI}(\nu_0)=7.42\times10^{-18}$,
      $\sigma_{\rm HeII}(\nu_0)=1.58\times10^{-18}\,{\rm cm^2}$.
      \flag{These are textbook values, not from the 57-paper corpus.}
\item \textbf{Thresholds}: $13.598$, $24.587$, $54.418\,$eV (NIST).
      \textbf{Cosmology}: $\Omega_{\rm b}h^2=0.02237$, $h=0.6736$,
      $Y_{\rm p}=0.2454$ \citep{Planck2018}.
\end{enumerate}

%=====================================================================
\section{Derivation}
%=====================================================================

The photoionization rate per target particle of species $i$ is
\begin{equation}
  \Gamma_i \;=\; \int_{\nu_i}^{\infty}
      \frac{4\pi J_\nu}{h\nu}\,\sigma_i(\nu)\,\mathrm{d}\nu ,
  \label{eq:gamma}
\end{equation}
with $J_\nu$ the mean specific intensity.  Writing $J_\nu=J_0 f(\nu)$ with
$f(\nu_{\rm HI})\equiv1$, the normalisation $J_0$ cancels in a ratio:
\begin{equation}
  \frac{\Gamma_i}{\GHI} \;=\;
  \frac{\displaystyle\int_{\nu_i}^{\infty} f(\nu)\,\sigma_i(\nu)\,\frac{\mathrm{d}\nu}{\nu}}
       {\displaystyle\int_{\nu_{\rm HI}}^{\infty} f(\nu)\,\sigma_{\rm HI}(\nu)\,\frac{\mathrm{d}\nu}{\nu}} .
  \label{eq:ratio}
\end{equation}
This is the whole point of the method: everything uncertain about the
\emph{amplitude} of the background drops out, and the answer depends only on
the shape $f$ and on atomic physics.  Absolute rates then follow from the
measured \GHI.

For the special case $f\propto\nu^{-\alpha}$ and $\sigma\propto\nu^{-3}$,
Eq.~\ref{eq:gamma} integrates in closed form to
\begin{equation}
  \Gamma_i \;=\; \frac{4\pi\,\sigma_i(\nu_i)\,J_\nu(\nu_i)}{h\,(\alpha+3)} ,
  \label{eq:closed}
\end{equation}
which is used only as a sanity limit; the numbers in
Section~\ref{sec:results} use the full cross-sections in Eq.~\ref{eq:ratio}.

%=====================================================================
\section{Verification}
%=====================================================================

\paragraph{Symbolic.}
Eq.~\ref{eq:closed} is verified with \texttt{sympy} by direct integration
(\texttt{simplify($\Gamma_{\rm sym}-$closed)$\,=\,$0}).

\paragraph{Dimensional.}
$[J_\nu]=\mathrm{erg\,cm^{-2}\,s^{-1}\,Hz^{-1}}=E/L^2$, $[h\nu]=E$,
$[\sigma]=L^2$, $[\mathrm{d}\nu]=1/T$, so
$[\Gamma]=(E/L^2)(1/E)(L^2)(1/T)=1/T$.  Confirmed with
\texttt{dimsys\_SI.equivalent\_dims}.

\paragraph{Independent numerical cross-check.}
\GHI\ is rebuilt from quantities that do \emph{not} include it, using the
local-source approximation $4\pi J_\nu=\epsilon_\nu^{\rm prop}\lambda^{\rm prop}$
with $\epsilon_{912}$ and $\lambda_{\rm mfp}$ from the same table
\citep{Gaikwad2023}:
\[
  \Gamma_{\rm HI}^{\rm rebuilt}=1.12\times10^{-13}\,{\rm s^{-1}}
  \qquad\text{vs}\qquad
  \Gamma_{\rm HI}^{\rm measured}=1.45\times10^{-13}\,{\rm s^{-1}},
\]
a ratio of $0.77$.  Agreement to $23\%$ is as good as this approximation can
be expected to do (it ignores the frequency dependence of the mean free path),
and it confirms that the emissivity, mean free path and photoionization rate in
that table are mutually consistent.

%=====================================================================
\section{Results}
\label{sec:results}
%=====================================================================

\begin{table}[htbp]
\centering
\caption{Photoionization rate ratios from Eq.~\ref{eq:ratio}.  Independent of
the background normalisation.}
\label{tab:ratios}
\begin{tabular}{@{}l c c@{}}
\toprule
Spectral basis & $\GHeI/\GHI$ & $\GHeII/\GHI$ \\
\midrule
(A) $\alpha_{\rm s}=1.4$ ($-1\sigma$)        & $0.706$ & $3.60\times10^{-2}$ \\
(A) $\alpha_{\rm s}=2.0$ (fiducial)          & $0.475$ & $1.57\times10^{-2}$ \\
(A) $\alpha_{\rm s}=2.6$ ($+1\sigma$)        & $0.322$ & $6.81\times10^{-3}$ \\
(B) BPASS stellar SED \citep{Eide2020}       & $0.651$ & $3.59\times10^{-3}$ \\
\bottomrule
\end{tabular}
\end{table}

\begin{table}[htbp]
\centering
\caption{Absolute photoionization rates, s$^{-1}$ per target particle.
\GHI\ is measured \citep{Gaikwad2023}; helium rates follow from
Table~\ref{tab:ratios}.  Quoted \GHI\ uncertainties are the published total
$1\sigma$; they scale the helium columns proportionally and are \emph{not}
the dominant error there --- the spectral basis is.}
\label{tab:abs}
\begin{tabular}{@{}l l c c@{}}
\toprule
$z$ & basis & $\GHeI$ (s$^{-1}$) & $\GHeII$ (s$^{-1}$) \\
\midrule
\multicolumn{4}{@{}l}{$\GHI = 1.45\,(+1.57/-0.87)\times10^{-13}\,{\rm s^{-1}}$ \hfill at $z=6.0$}\\
6.0 & $\alpha_{\rm s}=1.4$ & $1.02\times10^{-13}$ & $5.22\times10^{-15}$ \\
6.0 & $\alpha_{\rm s}=2.0$ & $6.89\times10^{-14}$ & $2.27\times10^{-15}$ \\
6.0 & $\alpha_{\rm s}=2.6$ & $4.67\times10^{-14}$ & $9.88\times10^{-16}$ \\
6.0 & BPASS                & $9.44\times10^{-14}$ & $5.21\times10^{-16}$ \\
\midrule
\multicolumn{4}{@{}l}{$\GHI = 5.01\,(+2.75/-2.32)\times10^{-13}\,{\rm s^{-1}}$ \hfill at $z=4.9$}\\
4.9 & $\alpha_{\rm s}=1.4$ & $3.54\times10^{-13}$ & $1.80\times10^{-14}$ \\
4.9 & $\alpha_{\rm s}=2.0$ & $2.38\times10^{-13}$ & $7.85\times10^{-15}$ \\
4.9 & $\alpha_{\rm s}=2.6$ & $1.62\times10^{-13}$ & $3.41\times10^{-15}$ \\
4.9 & BPASS                & $3.26\times10^{-13}$ & $1.80\times10^{-15}$ \\
\bottomrule
\end{tabular}
\end{table}

\paragraph{Photon-number rates.}
The literal ``ionizing photons per unit time'', for completeness.  At $z=6.0$
\citep{Gaikwad2023}, $\nion = 7.01\times10^{50}\,{\rm s^{-1}\,\cMpc^{-3}}$.
With $\langle n_{\rm H}\rangle = 1.895\times10^{-7}\,{\rm cm^{-3}}$ comoving
$=5.566\times10^{66}\,\cMpc^{-3}$ this is
$1.26\times10^{-16}\,{\rm s^{-1}}$ per hydrogen atom, i.e.\ $3.97$ ionizing
photons per H atom per Gyr.  The fraction emitted above each helium edge is
$30.6\%$ / $6.2\%$ for basis (A) at $\alpha_{\rm s}=2.0$, and
$32.5\%$ / $0.63\%$ for basis (B).

\paragraph{Pop III comparison (basis C).}
From \citet{Murphy2021}, Table~4 (non-rotating, $\alpha_{\rm ov}=0.1$):
$N_{\rm HeI}/N_{\rm H} = 0.384$ and $N_{\rm HeII}/N_{\rm H} = 0.0184$ for the
top-heavy SB13 IMF, and $0.310$ / $0.0096$ for a Salpeter slope.  A single
power law cannot reproduce both ratios: they imply an effective index of
$1.62$ across the He\,\textsc{i} edge but $2.88$ across the He\,\textsc{ii}
edge (SB13), confirming the curvature seen in basis (B).  Metal-free stars are
therefore harder than reionization-era populations in the He\,\textsc{i} band
but still show a He\,\textsc{ii} deficit.

%=====================================================================
\section{Final result}
%=====================================================================

At $z=6$, per target particle, in the ionized IGM:
\begin{align*}
  \GHI   &= 1.45\,(+1.57/-0.87)\times10^{-13}\ {\rm s^{-1}}
            &&\text{(measured)}\\
  \GHeI  &= (0.5\text{--}1.0)\times10^{-13}\ {\rm s^{-1}}
            &&\text{i.e. } (0.32\text{--}0.71)\,\GHI\\
  \GHeII &\lesssim (0.5\text{--}5)\times10^{-15}\ {\rm s^{-1}}
            &&\text{i.e. } (3.6\times10^{-3}\text{--}3.6\times10^{-2})\,\GHI
\end{align*}

\GHeI\ is robust: every basis agrees to within a factor $\sim2$, because the
$24.6$--$54.4\,$eV band is where stellar spectra are still strong and the IGM
is transparent.  \GHeII\ spans a factor $\sim10$ and its lowest value --- the
BPASS one --- is the physically appropriate one for a stellar population.

%=====================================================================
\section{Caveats that matter for the cosmic-ray comparison}
\label{sec:caveats}
%=====================================================================

\begin{enumerate}[leftmargin=1.7em]
\item \textbf{$\Gamma$ is per \emph{neutral} target, and so is $\zeta$.}  The
      comparison $\Gamma_i$ vs $\zeta_i$ is therefore apples-to-apples with no
      volume normalisation.  If a volumetric rate is wanted instead, multiply
      by the density of the \emph{neutral} species, not the total: inside
      ionized regions $x_{\rm HI}\sim10^{-4}$--$10^{-5}$, so
      $n_{\rm HI}\Gamma_{\rm HI}\ll n_{\rm H}\Gamma_{\rm HI}$.
\item \textbf{\GHeII\ is an upper bound.}  He\,\textsc{ii} is not reionized at
      $z\simeq6$ --- it completes only by $z\sim3$ \citep{Villasenor2021} ---
      so He\,\textsc{ii} is itself the dominant absorber above $54.4\,$eV and
      self-shields.  The optically-thin values in Table~\ref{tab:abs} therefore
      overestimate the true rate in the general IGM, by an amount that
      requires radiative transfer to quantify.  \flag{Not computed here.}
\item \textbf{These rates are zero in neutral regions.}  \GHI\ is a
      volume-average over \emph{ionized} gas.  Outside ionization fronts the
      photoionizing background collapses, while cosmic rays and X-rays --- with
      mean free paths far exceeding the bubble size --- do not.  This is the
      physically interesting regime for a cosmic-ray comparison: photons win by
      orders of magnitude inside H\,\textsc{ii} regions and are essentially
      absent outside them, whereas cosmic-ray ionization is comparatively
      uniform \citep{Sazonov2015,GesseyJones2023}.
\item \textbf{Helium abundance.}  $n_{\rm He}/n_{\rm H}=0.0813$ for
      $Y_{\rm p}=0.2454$, so even at equal $\Gamma$ the helium
      \emph{volumetric} ionization rate is $\sim12\times$ smaller than
      hydrogen's.
\end{enumerate}

\bibliographystyle{unsrtnat}
\bibliography{refs}
\end{document}
